\begin{problem}[Warmup]{(a)}
    Explain why the Constructible Number Theorem 38.3 guarantees that the following numbers are constructible: \[\sqrt{6}, \quad \sqrt[8]{6}, \quad \sqrt{2 + \sqrt[4]{5}}.\]
\end{problem}
\unsolved

\begin{problem}[Warmup]{(b)}
    What is the quadratic field extension $\mathbb{R}(i)$ usually known as?
\end{problem}
\unsolved

\begin{problem}[Warmup]{(c)}
    Does $\mathbb{C}$ admit any proper quadratic field extensions?
\end{problem}
\unsolved

\begin{problem}[Warmup]{(d)}
    Does $\mathbb{K}$ admit any proper quadratic field extensions? Does it admit any proper quadratic field extensions that are subfields of $\mathbb{R}$?
\end{problem}
\unsolved

\begin{problem}[Warmup]{(e)}
    Did we discuss a real number in this chapter that is \emph{not} constructible?
\end{problem}
\unsolved

\begin{problem}[Warmup]{(f)}
    Suppose that $\ell_1$ and $\ell_2$ are lines in the plane of a field $F$ of constructible numbers. What can you say about the intersection of $\ell_1$ and $\ell_2$?
\end{problem}
\unsolved

\begin{problem}[Warmup]{(g)}
    Suppose that $a$, $b$, and $c$ are constructible numbers, and $ax^2 + bx + c$ has real roots. Are these roots constructible numbers?
\end{problem}
\unsolved

\begin{problem}{1}
    Show that $\sqrt{5} \notin \mathbb{Q}(\sqrt{3})$ by showing that 5 cannot be the square
    of a number of the form $a + b\sqrt{3}$ where $a$ and $b$ are in $\mathbb{Q}$.
\end{problem}
\unsolved

\begin{problem}{2}
    If $F = \mathbb{Q}(\sqrt{3})$, describe the set of elements of the field $F(\sqrt{5})$.
    Show that this field is the same as the quadratic extension
    \[\mathbb{Q}(\sqrt{5})(\sqrt{3}).\]
\end{problem}
\unsolved

\begin{problem}{3}
    Generalize Exercise 2: Suppose that $p$ and $q$ are distinct positive prime integers;
    prove that
    \[\mathbb{Q}(\sqrt{p})(\sqrt{q}) = \mathbb{Q}(\sqrt{q})(\sqrt{p}),\]
    and describe the elements of the field.
\end{problem}
\unsolved

\begin{problem}{4}
    Give a sequence of numbers necessary to construct the number
    \[\sqrt[4]{2 + 4\sqrt{3}}.\]
    Give the corresponding sequence of fields.
\end{problem}
\unsolved

\begin{problem}{5}
    We've been able to give a nice description of the smallest field containing both
    $\sqrt{2}$ and $\mathbb{Q}$; namely, $\{a + b\sqrt{2} : a, b \in \mathbb{Q}\}$. In this
    problem we try to give a description of the smallest field containing $\sqrt[3]{2}$ and
    $\mathbb{Q}$. Why is $\{a + b\sqrt[3]{2} : a, b \in \mathbb{Q}\}$ not the answer? Now
    consider $\{a + b\sqrt[3]{2} + c\sqrt[3]{4} : a, b, c \in \mathbb{Q}\}$. When proving
    this is a field, the difficult part is showing that a typical element
    $a + b\sqrt[3]{2} + c\sqrt[3]{4}$, with $a, b, c \in \mathbb{Q}$ and not all zero, has a
    multiplicative inverse. You might not be successful in this part, and we will give an
    explicit hint below. But do at least verify that this set is indeed a commutative ring
    with unity.

    Here is a hint for how to show that a typical element indeed has a multiplicative
    inverse. In fact, the multiplicative inverse of $a + b\sqrt[3]{2} + c\sqrt[3]{4}$ is
    $d + e\sqrt[3]{2} + f\sqrt[3]{4}$ where
    \[d = \frac{-2bc + a^2}{2b^3 + 4c^3 + a^3 - 6bca}, \quad
      e = \frac{2c^2 - ba}{2b^3 + 4c^3 + a^3 - 6bca}, \quad\text{and}\quad
      f = \frac{ca - b^2}{2b^3 + 4c^3 + a^3 - 6bca}.\]
    You can verify this by computing
    $(a + b\sqrt[3]{2} + c\sqrt[3]{4})(d + e\sqrt[3]{2} + f\sqrt[3]{4})$ and seeing that it
    simplifies to 1. (This is obviously a non-trivial calculation!)

    We will show (in a less grubby manner) that this set is a field in Theorem 43.3. Note
    also that we have by no means claimed that $\sqrt[3]{2}$ is a constructible number (see
    Section 39.1).
\end{problem}
\unsolved

\begin{problem}{6}
    Describe the elements of the quadratic extension field $F(\sqrt[4]{2})$, where
    $F = \mathbb{Q}(\sqrt{2})$.
\end{problem}
\unsolved

\begin{problem}{7}
    Suppose that $a$, $b$, and $c$ are constructible numbers. Consider the polynomial
    $ax^4 + bx^2 + c$: Are its real roots constructible? (See Warm-up Exercise g above.)
\end{problem}
\unsolved

\begin{problem}{8}
    Argue that the linear equation that results when two circles are intersected
    algebraically gives the equation of the common chord of the two circles. (See the
    figure in the book.)
\end{problem}
\unsolved

\begin{problem}{9}
    Find the equation of tangent line to the circle
    \[{(x - 1)}^2 + y^2 = 5\]
    at the point $(3, 1)$ by algebra: This is the only line that intersects this circle in
    exactly this point. \emph{Note:} This was Rene Descartes' method for determining a
    tangent line. Doing this problem should make you appreciate the \emph{geometric}
    approach to finding the tangent (it is perpendicular to a radial line) and the
    \emph{calculus} approach (calculate $dy/dx$).
\end{problem}
\unsolved
